\section*{Step 255: Selberg Orthonormality and Cross-Correlation Extension}

\paragraph{Carrier.}
For primitive Selberg-class functions \(L_1,L_2\), define
\[
S_{12}(x)=\sum_{p\le x}
\frac{a_p(L_1)\overline{a_p(L_2)}}{p}.
\]
Full Selberg orthonormality consists of
\[
S_{LL}(x)=n_L\log\log x+O(1)
\]
on the diagonal and
\[
S_{12}(x)=O(1)
\]
for \(L_1\ne L_2\).  The constant \(n_L\) is the Selberg diagonal norm
constant, expected to be \(1\) for primitive functions; it is not the pole
order at \(s=1\).

\paragraph{Framework extension.}
The Selberg-Class Dichotomy Generalization from step 232 concerns one
\(L\)-function at a time: its RH analogue is target-equivalent within its own
family.  Selberg orthonormality is different.  It is a second-order carrier
whose residual is defined on pairs of primitive \(L\)-functions.

\paragraph{Candidate finding.}
\[
\textbf{Selberg-Class Cross-Correlation Extension.}
\]
Beyond per-\(L\) RH analogues, the Selberg class admits multi-\(L\) typed
residuals: coefficient orthogonality, full orthonormality, joint-zero
decorrelation, and related pairwise or multiway correlation conditions.

\paragraph{Evidence.}
The extension has two instances:
\[
\begin{array}{c|c|c}
\text{instance} & \text{step} & \text{residual}\\
\hline
\text{SOC off-diagonal} & 254 & S_{12}(x),\ L_1\ne L_2\\
\text{diagonal orthonormality} & 255 & S_{LL}(x)-n_L\log\log x
\end{array}
\]

\paragraph{Status.}
The finding is candidate, verified on two Selberg-class instances, and
corpus-pending.  It has been deposited in \texttt{anti_loc/findings_framework.md}.
No proof of SOC, orthonormality, GRH, or RH is claimed.
